\section{Experimental Protocol}
\label{app:reproducibility}
\begingroup
\small

This appendix records the protocol behind Section~\ref{sec:evidence} and the
evidence needed to audit an end-to-end system result. The current versioned
evaluation manifest covers GPQA-Diamond and LiveCodeBench; the HLE description
corresponds to the prior released scorecard retained in
Figure~\ref{fig:accuracy-scorecards}. The open reference implementation is available at
\url{https://github.com/vllm-project/semantic-router}. The evaluated
configurations were served through an OpenAI-compatible endpoint under the
shared alias \texttt{vllm-sr/auto}; each benchmark configuration nevertheless
has a distinct pool and topology and should be released under its own immutable
coordinate and digests.

\subsection{Constituent pools and trace budgets}

\paragraph{GPQA-Diamond.}
The benchmark is GPQA-Diamond~\citep{rein_2023_gpqa}.  The evaluated MoM
uses Gemini 3.1 Pro Preview, GPT-5.5, and Claude Opus 4.8.  Each request launches
one reasoning call to each constituent and then invokes Gemini for synthesis,
for four nominal model calls.  The output contract admits exactly one answer in
\{A,B,C,D\}.  The reported run contains 50 questions and is labeled
\emph{formal-50} to distinguish it from the complete 198-question benchmark.

\paragraph{LiveCodeBench.}
The evaluated window uses LiveCodeBench~\citep{jain_2024_livecodebench} v6,
January--April 2025, with 175 problems.  The pool again contains Gemini 3.1 Pro
Preview, GPT-5.5, and Claude Opus 4.8. Request signals select one of two
principal traces. The first uses three analysis calls followed by GPT-5.5
synthesis, for four calls in total; the second uses five breadth calls followed
by GPT-5.5 synthesis, for six calls in total. The contract requires one complete
Python solution with the problem's function or standard-input interface
preserved.  Sandboxed execution provides the task score.  Five missing
generations are assigned zero credit, so the denominator remains 175.

\paragraph{Humanity's Last Exam.}
The released scorecard reports a text-only filter over the 2,500-item release
of Humanity's Last Exam (HLE)~\citep{phan_2026_hle}, with 2,158 retained items.
The closed realization contains Claude Opus 4.8 and Gemini 3.1
Pro Preview. Depending on request risk, it obtains two, three, or four candidate
responses and then invokes Gemini for synthesis, for three to five nominal
calls; its scoring judge is also Gemini 3.1 Pro Preview. The hybrid realization
contains two independent GLM-5.2 sampling roles, GPT-5.5, Claude Opus 4.8, and
Gemini 3.1 Pro Preview. It obtains three or four candidate responses and then
invokes Gemini for synthesis, for four or five nominal calls.
The contract requires an explanation, exact answer, and calibrated confidence
field. The hybrid scoring path uses GPT-5.5. In both realizations, the judge
identity may also participate in the execution graph, so the reported numbers
are end-to-end system scores; comparisons intended to isolate constituent
quality should instead use a disjoint judge pool and report multi-judge
agreement.

For archival reproduction, each dataset must be pinned by repository revision,
configuration, filter predicate, ordered item identifiers, and an identifier
digest. In particular, \emph{LiveCodeBench v6} here denotes the 175-item
January--April 2025 configuration rather than the cumulative
\texttt{release\_v6}; the HLE count is produced by the text-only filter rather
than a benchmark-defined split; and the GPQA formal-50 item list requires an
explicit selection record. These public benchmarks predate the evaluated 2026
constituents, so no result is interpreted as contamination-free
generalization.

All constituent calls in these three experiments use the same managed model
gateway.  Thus the experiments test heterogeneous model identities and trace
topologies, while cross-provider and cross-hardware portability remain separate
evaluation dimensions.

\subsection{Request, trace, and scoring records}

A reproducible and independently auditable MoM evaluation requires more than
final predictions. Each
request record should contain:

\begin{itemize}
  \item dataset version, split, immutable example identifier, and input digest;
  \item logical MoM identity, behavior variant, request preference, bundle digest, and deployment-lock
  digest;
  \item signals, projection values, matched decision, selected topology, and
  hard-constraint outcomes;
  \item every constituent model and service revision, prompt digest, sampling
  parameters, retry index, stop reason, and terminal status;
  \item raw response or policy-compliant encrypted/redacted equivalent, parsed
  answer, contract-validation result, repair, fallback, and abstention events;
  \item input, cached-input, reasoning, and output token counts when available;
  \item per-call and end-to-end latency, timeouts, rate limits, and failures;
  \item realized monetary and, where estimable, energy cost under the declared
  accounting model;
  and
  \item scorer and judge versions, raw score, aggregation rule, and uncertainty.
\end{itemize}

Run-level metadata should additionally record engine and dependency versions,
model inventory, start and end times, concurrency, seeds, counts of attempted,
completed, and missing items, and checksums for every output shard. A run is complete
only when the expected item identifiers and denominator match; the presence of
an aggregate score file alone is insufficient.

\subsection{Routing-aware fleet-planning protocol}

The capacity example uses the empirical request-length distribution from
LMSYS-Chat-1M, collected partly through Chatbot
Arena~\citep{zheng_2023_lmsyschat,chiang_2024_chatbotarena}, at an arrival rate
of 100 requests per second and a 500 ms time-to-first-token target. The modeled
worker serves a 70B model with eight-way tensor parallelism on A100 GPUs. The
homogeneous baseline assigns every request to a 65K-context pool.  The routed
plan assigns short requests to a 4K-context pool and long requests to the 65K
pool, then sweeps the threshold subject to the same latency target.  GPU cost is
fixed at \$2.21 per hour.  Both reported plans are continuous analytical
capacity estimates before rounding to whole eight-GPU workers; they are not
deployable cluster counts or measurements from a production cluster.

For deployment use, the planner should first be calibrated against observed
prefill/decode throughput and tail latency, then validated on held-out traces.
A complete study should repeat the sweep over arrival processes, context
distributions, SLOs, accelerator types, model profiles, failures, and price
assumptions, and should compare analytical predictions with discrete-event and
live-cluster measurements.

\subsection{Release contract}

An independently checkable release should publish the canonical MoM bundle,
deployment binding and lock, constituent inventory, prompts, contracts,
benchmark manifests, per-item traces, and deterministic report generator.
Credentials, authorization headers, and private request content must never be
embedded; secret references and policy-preserving redaction retain the logical
identity needed for verification.  All tables and figures should be regenerated
from immutable request and run records, and their reported digests should match
the released artifacts.
\endgroup
